\chapter*{Conclusion}
\phantomsection
\addcontentsline{toc}{chapter}{Conclusion}
\markboth{Conclusion}{}

The capability this project set out to reach was already present in the hardware. Instruction trace is implemented across a large part of the STM32 range, and on the families carrying a Trace Memory Controller a capture may be taken without leaving the chip. What was missing was a way to reach it with the tools and the probe an ordinary debug session already uses.

The work produced three results. An application note assembles into one document the procedure for using instruction trace on these devices, which the architecture specifications and the reference manuals state separately. A proof of concept demonstrates that a capture may be obtained from an on-chip sink, decoded against the program image, and reconstructed into the instructions that executed. A debugger carries that capability into a development tool, alongside breakpoints, variable inspection, memory and peripheral views, and disassembly.

Section \ref{sec:testing} states what the result achieves. A capture of 128 bytes taken over the ordinary debug connection reconstructs 114 executed instructions, each attributed to a function, a source file and a line. No trace pins are used, and no probe beyond the one embedded on the development board.

Two conclusions carry beyond the project. The first is that the barrier to instruction trace on these devices was never the hardware. It was the distance between what the documentation states and what a developer must do, together with the absence of an open implementation to close that distance. The second is that reconstructing execution is not a decoding problem alone. The trace reports which branches were taken and nothing further, so the program image supplies everything else, and the result is only ever as good as the correspondence between the two.

What remains is stated where it belongs. The pipeline captures without filtering, discards the timing elements it decodes, and reads the trace topology from a configuration file, where the device's own ROM tables could supply it. Each of these is a boundary of the implementation and not of the approach, and Section \ref{subsec:limits} states what closing them would require.

Those boundaries also mark the directions the work opens. Retaining the timing elements would bring profiling from the same captures. Presenting exceptions explicitly would extend the reconstruction across interrupt entry. Reading the component topology from the device would remove the last per-device configuration. Each of these builds on the pipeline this project established, and none of them revisits the barrier the project set out to remove, since that barrier is now gone.
